\documentclass[11pt]{article}
\usepackage[a4paper,margin=1in]{geometry}
\usepackage{amsmath,amssymb,amsthm,mathtools}
\usepackage{graphicx}
\usepackage{microtype}
\usepackage[hidelinks]{hyperref}

\newtheorem{theorem}{Theorem}
\newtheorem{corollary}[theorem]{Corollary}
\newcommand{\Inf}{\operatorname{Inf}}
\newcommand{\tr}{\operatorname{tr}}

\title{Quantum KKL for Anticommuting-Pauli Observables\\
\large A substantial partial result for the Montanaro--Osborne conjecture}
\author{Automated Mathematics Run \texttt{fa\_banach\_001}, lane 2}
\date{9 August 2026}

\begin{document}
\maketitle

\begin{center}
\fbox{\parbox{0.9\linewidth}{\textbf{Status.}
Candidate partial result, likely valid. The full quantum $L^2$-KKL
conjecture remains open.}}
\end{center}

\section{Source problem}

Rouz\'e, Wirth, and Zhang study quantum Talagrand, KKL, and Friedgut
inequalities in arXiv:2209.07279, later published in \emph{Communications in
Mathematical Physics} \cite{RWZ}. Their Remark 3.10 on page 13 recalls
Montanaro and Osborne's Conjecture 3 from Section 12 of
\cite{MO}. The source passage is reproduced below.

\begin{figure}[ht]
  \centering
  \includegraphics[width=\linewidth]{figures/open_problem_crop.png}
  \caption{Full-width crop of Remark 3.10 on page 13 of the source paper.}
\end{figure}

\paragraph{Transcription.}
For a balanced quantum Boolean function
$A\in M_2(\mathbb C)^{\otimes n}$, meaning $A=A^*$, $A^2=I$, and
$\tr(A)=0$, is there a universal constant $C>0$ such that
\[
  \max_{1\le j\le n}\Inf_j^2(A)
  \ge C\frac{\log n}{n}?
\]
Here $d_j$ deletes the Pauli--Fourier terms that are identity on qubit $j$,
and $\Inf_j^2(A)=\lVert d_jA\rVert_2^2$ for normalized Schatten norm.
The source proves the corresponding result for $L^p$-influences when
$1\le p<2$ and explicitly states that this $p=2$ endpoint remains open
\cite{RWZ}. A later random-restriction treatment also states that the exact
conjecture remains open \cite{JLLZ}.

\section{Definitions}

Let $\mathcal P_n=\{I,X,Y,Z\}^{\otimes n}$ be the Hermitian Pauli strings.
With normalized trace $2^{-n}\tr$, this is an orthonormal basis. For
$\Gamma\in\mathcal P_n$, write
\[
 \operatorname{supp}(\Gamma)=\{j:\Gamma_j\ne I\}.
\]
For a Pauli expansion $A=\sum_\Gamma\widehat A(\Gamma)\Gamma$,
orthogonality gives
\begin{equation}\label{eq:influence-fourier}
 \Inf_j^2(A)
 =\sum_{\Gamma:\,j\in\operatorname{supp}(\Gamma)}
   |\widehat A(\Gamma)|^2.
\end{equation}

We call $A$ an \emph{anticommuting-Pauli observable} if
\begin{equation}\label{eq:A}
 A=\sum_{s=1}^m c_s\Gamma_s,
 \qquad c_s\in\mathbb R,
 \qquad \sum_{s=1}^m c_s^2=1,
\end{equation}
where the $\Gamma_s$ are distinct nonidentity Pauli strings satisfying
$\Gamma_s\Gamma_t=-\Gamma_t\Gamma_s$ whenever $s\ne t$.
Such observables arise naturally from Clifford generators and
fermion-to-qubit encodings.

\section{Proof intuition}

The unitarity of \eqref{eq:A} is immediate: all cross terms cancel in
$A^2$. The influence problem becomes a weighted incidence problem. Give
$\Gamma_s$ weight $w_s=c_s^2$. Every pair of strings must locally
anticommute on at least one qubit. At a fixed qubit, pairs locally
anticommute only across the three Pauli labels $X,Y,Z$; if their three total
weights are $x,y,z$, the covered pair weight is
$xy+xz+yz\le(x+y+z)^2/3$. Thus $n$ qubits can cover all pairs only if some
qubit carries substantial total weight. A single dominant coefficient is
handled separately by observing that one qubit in its support already has at
least that much influence. Balancing these two alternatives yields an
explicit $\Theta(n^{-1/2})$ bound.

The argument does not extend automatically to a general quantum Boolean
function. In a general Pauli expansion, commuting pairs may have identical
products with cancelling signed coefficients in $A^2=I$. Controlling that
cancellation geometry is precisely what remains open.

\section{The partial theorem}

\begin{theorem}[Anticommuting-Pauli quantum KKL]\label{thm:main}
Let $A$ be an anticommuting-Pauli observable on $n$ qubits. Then $A$ is a
balanced quantum Boolean function and
\begin{equation}\label{eq:rho}
 \max_{1\le j\le n}\Inf_j^2(A)
 \ge \rho_n
 :=\frac{\sqrt{24n+9}-3}{4n}
 =\frac{6}{\sqrt{24n+9}+3}.
\end{equation}
Consequently,
\begin{equation}\label{eq:root}
 \max_j\Inf_j^2(A)
 \ge \frac{6}{\sqrt{33}+3}\frac1{\sqrt n}
 \ge \frac{3e}{\sqrt{33}+3}\frac{\log n}{n}.
\end{equation}
In particular, the Montanaro--Osborne conjecture holds for this class, with
the stronger order $n^{-1/2}$.
\end{theorem}

\begin{proof}
Hermiticity is clear. Since $\Gamma_s^2=I$ and distinct terms anticommute,
\[
 A^2=\sum_s c_s^2I+
 \sum_{s<t}c_sc_t(\Gamma_s\Gamma_t+\Gamma_t\Gamma_s)=I.
\]
Every $\Gamma_s$ is nonidentity and hence traceless, so $\tr(A)=0$.

Put $w_s=c_s^2$, so $w_s\ge0$ and $\sum_sw_s=1$. For each qubit $j$ and
$\alpha\in\{X,Y,Z\}$, define
\[
 u_{j,\alpha}=\sum_{s:(\Gamma_s)_j=\alpha}w_s,
 \qquad
 q_j=u_{j,X}+u_{j,Y}+u_{j,Z}.
\]
By \eqref{eq:influence-fourier}, $q_j=\Inf_j^2(A)$. Set
$Q=\max_jq_j$.

Two Pauli strings anticommute exactly when the number of coordinates at
which they carry two distinct nonidentity Pauli labels is odd. Therefore
each unordered pair $s<t$ is locally anticommuting on at least one qubit.
Counting the pair weight $w_sw_t$ at every such qubit gives
\begin{align}
 \sum_{s<t}w_sw_t
 &\le \sum_{j=1}^n
 \bigl(u_{j,X}u_{j,Y}+u_{j,X}u_{j,Z}+u_{j,Y}u_{j,Z}\bigr)\notag\\
 &\le \frac13\sum_{j=1}^nq_j^2
 \le \frac n3Q^2.\label{eq:cover}
\end{align}
The middle inequality follows from
$xy+xz+yz=( (x+y+z)^2-(x^2+y^2+z^2))/2\le(x+y+z)^2/3$.

Let $a=\max_sw_s$. Since a string attaining $a$ is nonidentity, some
qubit lies in its support, and hence $Q\ge a$. Also
$\sum_sw_s^2\le a\sum_sw_s=a$. Thus \eqref{eq:cover} and
\[
 \sum_{s<t}w_sw_t=\frac{1-\sum_sw_s^2}{2}
\]
imply
\[
 Q\ge \max\left\{a,\sqrt{\frac{3(1-a)}{2n}}\right\}.
\]
The first branch increases in $a$ and the second decreases. Their intersection
solves
$a^2=3(1-a)/(2n)$, whose positive root is
\[
 a=\frac{\sqrt{24n+9}-3}{4n}
  =\frac{6}{\sqrt{24n+9}+3}=\rho_n.
\]
This proves \eqref{eq:rho}.

For $n\ge1$, $\sqrt{24n+9}+3\le(\sqrt{33}+3)\sqrt n$, giving the first
inequality in \eqref{eq:root}. Finally, the maximum of
$(\log x)/\sqrt x$ on $x\ge1$ is $2/e$, so
$n^{-1/2}\ge e(\log n)/(2n)$. This proves the second inequality.
\end{proof}

\begin{corollary}[Equal-coefficient incidence form]
If $\Gamma_1,\dots,\Gamma_m$ are distinct nonidentity pairwise
anticommuting Pauli strings and
$A=m^{-1/2}\sum_{s=1}^m\Gamma_s$, then some qubit belongs to at least
$m\rho_n$ of their supports.
\end{corollary}

\begin{proof}
For this $A$, equation \eqref{eq:influence-fourier} says that the influence
of qubit $j$ is exactly the fraction of strings whose support contains $j$.
Apply Theorem~\ref{thm:main}.
\end{proof}

\section{A current construction stress test}

Slote's August 2026 Parseval-limit construction \cite{Slote} produces
homogeneous unitary Pauli expansions with exponentially many equal
coefficients. It begins with an equal-weight, pairwise-anticommuting core
$P_1,\dots,P_L$, then attaches many commuting diagonal suffix terms. This
makes it a natural possible counterexample mechanism.

It does not evade the theorem. Each core index $\ell$ carries total squared
Fourier mass $1/L$, so the influence of a core qubit is exactly its weighted
incidence among the $P_\ell$. Applying Theorem~\ref{thm:main} to the
normalized core gives a core influence of order at least the inverse square
root of the number of core qubits. This is much larger than the KKL scale for
the full construction. Thus the construction refutes quantum Fourier
entropy--influence but not quantum KKL.

\section{Verification, novelty, and limitations}

The script \texttt{code/verify\_small\_cases.py} exhausts every subset of the
15 nonidentity two-qubit Pauli strings, retains the pairwise-anticommuting
families, and checks the theorem and the intermediate pair-cover inequality
for deterministic weight samples. It also checks random greedy families on
three qubits. No violation was found. This is a finite consistency check and
is not used in the proof.

A bounded novelty search was run through 9 August 2026. The sources checked
included arXiv:0810.2435, 2209.07279, 2411.12399, and 2608.01424, plus exact
and close web searches for ``pairwise anticommuting Pauli influence KKL'',
``anticommuting Pauli quantum KKL'', and ``Clifford quantum KKL influence''.
The searches found standard uses of anticommuting Pauli families and the
known $p<2$ quantum KKL variants, but not the weighted incidence theorem
above. Novelty confidence is moderate, not definitive.

The full $L^2$ conjecture remains open. The theorem covers only Pauli
expansions whose entire nonzero support is pairwise anticommuting. It does not
control general unitaries with cancellations among products of commuting
Pauli terms. The exponent $1/2$ is a universal lower bound for this class; no
claim of sharpness of the exact constant is made.

\paragraph{Human-review recommendation.}
Review as a likely-valid substantial partial result. The highest-value checks
are the inequality direction in the weighted pair cover \eqref{eq:cover}, the
Fourier normalization in \eqref{eq:influence-fourier}, and the one-variable
optimization producing \(\rho_n\).

\begin{thebibliography}{9}
\bibitem{MO}
A.~Montanaro and T.~J. Osborne,
\emph{Quantum Boolean functions}, Chicago Journal of Theoretical Computer
Science (2010), arXiv:0810.2435.

\bibitem{RWZ}
C.~Rouz\'e, M.~Wirth, and H.~Zhang,
\emph{Quantum Talagrand, KKL and Friedgut's theorems and the learnability of
quantum Boolean functions}, Communications in Mathematical Physics
\textbf{405} (2024), Article 95, arXiv:2209.07279.

\bibitem{JLLZ}
Y.~Jiao, W.~Lin, S.~Luo, and D.~Zhou,
\emph{Quantum KKL-type inequalities revisited}, arXiv:2411.12399 (2024).

\bibitem{Slote}
J.~Slote,
\emph{Dense Hamiltonians at the Parseval limit: The noncommutative BH
constant is exponential and the quantum FEI conjecture is false},
arXiv:2608.01424 (2026).
\end{thebibliography}

\end{document}
